\documentclass[11pt,a4paper]{report}

\usepackage[margin=1in]{geometry}
\usepackage{amsmath,amssymb}
\usepackage{booktabs}
\usepackage{listings}
\usepackage{xcolor}
\usepackage{hyperref}
\usepackage{fancyhdr}
\usepackage{titlesec}
\usepackage{microtype}

% Color definitions for syntax highlighting
\definecolor{codegreen}{rgb}{0.1,0.6,0.2}
\definecolor{codegray}{rgb}{0.5,0.5,0.5}
\definecolor{codepurple}{rgb}{0.58,0,0.82}
\definecolor{backcolour}{rgb}{0.97,0.97,0.98}

\lstdefinestyle{sqlstyle}{
    backgroundcolor=\color{backcolour},   
    commentstyle=\color{codegreen},
    keywordstyle=\color{blue}\bfseries,
    numberstyle=\tiny\color{codegray},
    stringstyle=\color{codepurple},
    basicstyle=\ttfamily\footnotesize,
    breakatwhitespace=false,         
    breaklines=true,                 
    captionpos=b,                    
    keepspaces=true,                 
    numbers=left,                    
    numbersep=5pt,                  
    showspaces=false,                
    showstringspaces=false,
    showtabs=false,                  
    tabsize=2,
    frame=lines
}

\lstset{style=sqlstyle}

\pagestyle{fancy}
\fancyhf{}
\rhead{SmartFlow Grid: DBMS Case Study}
\lhead{Abraham \& Kundu}
\cfoot{\thepage}

\begin{document}

% -----------------------------------------------------------------------------
% Title Page
% -----------------------------------------------------------------------------
\begin{titlepage}
    \centering
    \vspace*{1.5cm}
    
    {\Huge\bfseries SmartFlow Grid:\\[0.3cm] Intelligent IoT Traffic Management System \par}
    \vspace{0.8cm}
    {\Large\itshape Advanced Relational Database Architecture, Real-Time Telemetry Ingestion, and Automated PL/SQL Signal Actuation \par}
    
    \vspace{2cm}
    {\large\textbf{Academic Case Study Comprehensive Submission (Reviews 1, 2, and 3)}\par}
    \vspace{0.3cm}
    {\large Course: Database Management Systems \par}
    
    \vspace{2.5cm}
    \begin{minipage}{0.8\textwidth}
        \centering
        \textbf{Prepared by:}\\[0.4cm]
        {\large\textbf{Taanush Emmanuel Abraham}}\\[0.1cm]
        Registration Number: \texttt{24BCE0708}\\[0.6cm]
        {\large\textbf{Tanishka Kundu}}\\[0.1cm]
        Registration Number: \texttt{24BCE0730}
    \end{minipage}
    
    \vfill
    {\large October 2026 \par}
\end{titlepage}

% -----------------------------------------------------------------------------
% Abstract
% -----------------------------------------------------------------------------
\begin{abstract}
Urban traffic congestion represents a multi-billion dollar economic drain and public safety liability, driven largely by antiquated fixed-time traffic light cycles that fail to adjust to dynamic, non-stationary vehicular flows. The \textbf{SmartFlow Grid} provides an advanced relational database system designed to model physical road intersection topologies, ingest real-time IoT sensor telemetry, and execute real-time automated signal actuation through server-side procedural logic. This project delivers a production-grade dual implementation across \textbf{Oracle Database (PL/SQL)} and \textbf{PostgreSQL (PL/pgSQL)}. The system enforces strict integrity constraints, models dedicated directional lanes and signal phases, computes weighted moving average (WMA) congestion indices, facilitates automated emergency corridor clearance, conducts cursor-driven regional zone surveillance, and maintains an asynchronous shadow audit log. Additionally, this report establishes a formal mathematical normalization proof from an unnormalized Universal Relation $R$ through 1NF, 2NF, 3NF, and Boyce-Codd Normal Form (BCNF).
\end{abstract}

\tableofcontents
\newpage

% -----------------------------------------------------------------------------
% Chapter 1
% -----------------------------------------------------------------------------
\chapter{Domain Analysis \& Problem Formulation}

\section{The Urban Traffic Gridlock Problem}
Conventional municipal traffic management systems operate on static time-of-day interval schedules configured around historical averages. These open-loop controllers suffer from three fundamental systemic flaws:
\begin{enumerate}
    \item \textbf{Asymmetric Volume Spikes:} Unanticipated events such as accidents, disabled vehicles, or construction bottlenecks create severe queue backups on specific approaches, while opposing green phases remain underutilized.
    \item \textbf{Emergency Preemption Delays:} First responders (ambulances, fire engines, law enforcement) lose critical minutes traversing congested corridors without automated, fail-safe signal override preemption.
    \item \textbf{Excess Emissions and Fuel Consumption:} Prolonged idling at uncoordinated signals accounts for substantial greenhouse gas emissions and economic productivity losses.
\end{enumerate}

\section{The SmartFlow Grid Architectural Paradigm}
The SmartFlow Grid addresses these challenges by transforming every physical intersection into a monitored relational node. Directional approach lanes are equipped with IoT inductive loops and video sensors that stream vehicular count telemetry into a central database. By housing business and actuation logic directly inside the database via PL/SQL procedures, functions, explicit cursors, and triggers, the system achieves sub-millisecond actuation response times and strict transaction atomicity without reliance on external application layers.

% -----------------------------------------------------------------------------
% Chapter 2
% -----------------------------------------------------------------------------
\chapter{Requirements Analysis \& Business Rules}

\section{Functional Requirements}
\begin{itemize}
    \item \textbf{FR1 (Infrastructure Modeling):} Model junctions (\texttt{Intersection}), surveillance zones (\texttt{Grid\_Zone}), and directional lanes (\texttt{Lane}).
    \item \textbf{FR2 (Lane-Level Phase Actuation):} Maintain dedicated 1:1 signal phase controllers (\texttt{Signal\_Phase}) per lane.
    \item \textbf{FR3 (High-Frequency Telemetry Ingestion):} Ingest continuous time-series logs (\texttt{Sensor\_Log}) linked to specific lanes.
    \item \textbf{FR4 (Congestion Classification):} Calculate a weighted moving average (WMA) across recent sensor readings to categorize congestion into \texttt{'Low'}, \texttt{'Medium'}, and \texttt{'High'}.
    \item \textbf{FR5 (Emergency Corridor Override):} Provide an automated procedure (\texttt{proc\_Emergency\_Override}) that forces conflicting lanes to \texttt{'RED'} while clearing emergency paths to \texttt{'GREEN'}.
    \item \textbf{FR6 (Regional Zone Surveillance):} Use explicit cursors (\texttt{proc\_Check\_Zone\_Diversion}) to evaluate cumulative zone saturation and emit automated diversion recommendations.
    \item \textbf{FR7 (Immutable Shadow Audit):} Deploy a reactive trigger (\texttt{trg\_Signal\_Phase\_Audit}) to record all signal state mutations into \texttt{Audit\_Traffic\_Log}.
\end{itemize}

\section{Integrity Constraints \& Validation Rules}
\begin{itemize}
    \item \textbf{Primary Keys:} Enforced on all tables (\texttt{Intersection\_ID}, \texttt{Lane\_ID}, \texttt{Phase\_ID}, \texttt{Log\_ID}, \texttt{Audit\_ID}).
    \item \textbf{Foreign Keys:} Cascading deletions on dependent tables; 1:1 unique constraint on \texttt{Signal\_Phase(Lane\_ID)}.
    \item \textbf{Check Constraints:}
    \begin{itemize}
        \item $\texttt{Current\_State} \in \{\text{'RED'}, \text{'YELLOW'}, \text{'GREEN'}\}$
        \item $\texttt{Direction} \in \{\text{'NORTHBOUND'}, \text{'SOUTHBOUND'}, \text{'EASTBOUND'}, \text{'WESTBOUND'}\}$
        \item $\texttt{Capacity\_Limit} > 0$, $\texttt{Car\_Count} \ge 0$, $\texttt{Duration\_Seconds} > 0$
        \item $\texttt{Is\_Emergency\_Corridor} \in \{0, 1\}$
    \end{itemize}
\end{itemize}

% -----------------------------------------------------------------------------
% Chapter 3
% -----------------------------------------------------------------------------
\chapter{Database Design \& EER Modeling}

\section{Relational Schema Overview}
The system architecture comprises five core relational entities structured to eliminate redundancy while capturing high-throughput operational dynamics:
\begin{enumerate}
    \item $\texttt{Intersection}(\underline{\text{Intersection\_ID}}, \text{Location\_Name}, \text{Grid\_Zone})$
    \item $\texttt{Lane}(\underline{\text{Lane\_ID}}, \text{Intersection\_ID}, \text{Direction}, \text{Capacity\_Limit}, \text{Is\_Emergency\_Corridor})$
    \item $\texttt{Signal\_Phase}(\underline{\text{Phase\_ID}}, \text{Lane\_ID}, \text{Intersection\_ID}, \text{Current\_State}, \text{Duration\_Seconds}, \text{Last\_Updated})$
    \item $\texttt{Sensor\_Log}(\underline{\text{Log\_ID}}, \text{Lane\_ID}, \text{Car\_Count}, \text{Recorded\_At})$
    \item $\texttt{Audit\_Traffic\_Log}(\underline{\text{Audit\_ID}}, \text{Phase\_ID}, \text{Old\_State}, \text{New\_State}, \text{Changed\_At}, \text{Trigger\_Source})$
\end{enumerate}

% -----------------------------------------------------------------------------
% Chapter 4
% -----------------------------------------------------------------------------
\chapter{Formal Normalization Proofs (1NF to BCNF)}

\section{The Universal Relation ($R$)}
We define the unnormalized flat relation $R$:
\begin{align*}
R = \{ & \text{Intersection\_ID}, \text{Location\_Name}, \text{Grid\_Zone}, \text{Lane\_ID}, \text{Direction}, \text{Capacity\_Limit}, \\
       & \text{Is\_Emergency}, \text{Phase\_ID}, \text{Current\_State}, \text{Duration}, \text{Last\_Updated}, \text{Log\_ID}, \\
       & \text{Car\_Count}, \text{Recorded\_At} \}
\end{align*}

\section{Set of Functional Dependencies ($F$)}
Based on the physical domain constraints, $F$ comprises:
\begin{align}
FD_1 &: \text{Intersection\_ID} \rightarrow \text{Location\_Name}, \text{Grid\_Zone} \\
FD_2 &: \text{Location\_Name} \rightarrow \text{Intersection\_ID}, \text{Grid\_Zone} \\
FD_3 &: \text{Lane\_ID} \rightarrow \text{Intersection\_ID}, \text{Direction}, \text{Capacity\_Limit}, \text{Is\_Emergency} \\
FD_4 &: \text{Phase\_ID} \rightarrow \text{Lane\_ID}, \text{Intersection\_ID}, \text{Current\_State}, \text{Duration}, \text{Last\_Updated} \\
FD_5 &: \text{Lane\_ID} \rightarrow \text{Phase\_ID} \\
FD_6 &: \text{Log\_ID} \rightarrow \text{Lane\_ID}, \text{Car\_Count}, \text{Recorded\_At} \\
FD_7 &: \{\text{Lane\_ID}, \text{Recorded\_At}\} \rightarrow \text{Car\_Count}, \text{Log\_ID}
\end{align}

\section{Candidate Key Analysis}
Computing the attribute closure of candidate determinant $\{\text{Log\_ID}\}$:
$$\{\text{Log\_ID}\}^+ = R$$
Thus, the primary candidate key is $K_1 = \{\text{Log\_ID}\}$. Alternate candidate keys include $K_2 = \{\text{Lane\_ID}, \text{Recorded\_At}\}$ and $K_3 = \{\text{Phase\_ID}, \text{Recorded\_At}\}$.

\section{First Normal Form (1NF)}
All attribute domains consist strictly of atomic values without repeating groups or nested arrays. Hence, $R$ satisfies \textbf{1NF}.

\section{Second Normal Form (2NF)}
Under composite candidate key $K_2 = \{\text{Lane\_ID}, \text{Recorded\_At}\}$, $FD_3$ indicates:
$$\text{Lane\_ID} \rightarrow \text{Capacity\_Limit}$$
Here, $\text{Capacity\_Limit}$ is non-prime and depends on a proper subset of $K_2$, violating 2NF (Partial Dependency). Decomposing $R$ yields:
\begin{itemize}
    \item $R_{\text{Telemetry}}(\underline{\text{Log\_ID}}, \text{Lane\_ID}, \text{Car\_Count}, \text{Recorded\_At})$
    \item $R_{\text{Infra}}(\underline{\text{Lane\_ID}}, \text{Intersection\_ID}, \text{Location\_Name}, \text{Grid\_Zone}, \text{Direction}, \text{Capacity\_Limit}, \dots)$
\end{itemize}
This decomposition eliminates partial dependencies, satisfying \textbf{2NF}.

\section{Third Normal Form (3NF)}
In $R_{\text{Infra}}$, candidate key is $\text{Lane\_ID}$. However, $FD_1$ shows:
$$\text{Intersection\_ID} \rightarrow \text{Location\_Name}, \text{Grid\_Zone}$$
$\text{Intersection\_ID}$ is not a superkey of $R_{\text{Infra}}$, creating a Transitive Dependency ($\text{Lane\_ID} \rightarrow \text{Intersection\_ID} \rightarrow \text{Location\_Name}$). Decomposing into separate entities:
\begin{enumerate}
    \item $\texttt{Intersection}(\underline{\text{Intersection\_ID}}, \text{Location\_Name}, \text{Grid\_Zone})$
    \item $\texttt{Lane}(\underline{\text{Lane\_ID}}, \text{Intersection\_ID}, \text{Direction}, \text{Capacity\_Limit}, \text{Is\_Emergency\_Corridor})$
    \item $\texttt{Signal\_Phase}(\underline{\text{Phase\_ID}}, \text{Lane\_ID}, \text{Intersection\_ID}, \text{Current\_State}, \text{Duration}, \text{Last\_Updated})$
    \item $\texttt{Sensor\_Log}(\underline{\text{Log\_ID}}, \text{Lane\_ID}, \text{Car\_Count}, \text{Recorded\_At})$
\end{enumerate}
In every dependency $X \rightarrow A$, $X$ is a superkey. Thus, the schema satisfies \textbf{3NF}.

\section{Boyce-Codd Normal Form (BCNF)}
In each decomposed relation, every determinant is a superkey:
\begin{itemize}
    \item In $\texttt{Intersection}$, keys are $\text{Intersection\_ID}$ and $\text{Location\_Name}$. Both are superkeys.
    \item In $\texttt{Lane}$, determinant $\text{Lane\_ID}$ is a superkey.
    \item In $\texttt{Signal\_Phase}$, determinants $\text{Phase\_ID}$ and $\text{Lane\_ID}$ are both superkeys.
    \item In $\texttt{Sensor\_Log}$, determinants $\text{Log\_ID}$ and $\{\text{Lane\_ID}, \text{Recorded\_At}\}$ are superkeys.
\end{itemize}
Therefore, the schema is strictly in \textbf{BCNF}.

\section{Lossless Join and Dependency Preservation Proofs}
\begin{itemize}
    \item \textbf{Lossless Join:} For each decomposition pair $(R_1, R_2)$, $(R_1 \cap R_2) \rightarrow R_1$ or $(R_1 \cap R_2) \rightarrow R_2$. By Heath's Theorem, the decomposition is guaranteed lossless.
    \item \textbf{Dependency Preservation:} All functional dependencies in $F$ are preserved within individual decomposed tables without requiring inter-table joins.
\end{itemize}

% -----------------------------------------------------------------------------
% Chapter 5
% -----------------------------------------------------------------------------
\chapter{Implementation Code Catalog}

\section{Live Dashboard View (\texttt{vw\_Dashboard\_Live})}
\begin{lstlisting}[language=SQL]
CREATE OR REPLACE VIEW vw_Dashboard_Live AS
SELECT 
    i.Intersection_ID,
    i.Location_Name,
    l.Lane_ID,
    l.Direction,
    sp.Current_State AS Active_Color,
    sp.Duration_Seconds,
    s.Car_Count AS Current_Volume,
    ROUND((s.Car_Count / l.Capacity_Limit) * 100, 2) AS Saturation_Pct
FROM Signal_Phase sp
JOIN Lane l ON sp.Lane_ID = l.Lane_ID
JOIN Intersection i ON sp.Intersection_ID = i.Intersection_ID
JOIN Sensor_Log s ON l.Lane_ID = s.Lane_ID
WHERE sp.Current_State = 'GREEN'
  AND s.Recorded_At = (
      SELECT MAX(s2.Recorded_At) 
      FROM Sensor_Log s2 WHERE s2.Lane_ID = l.Lane_ID
  );
\end{lstlisting}

\section{Emergency Override Stored Procedure}
\begin{lstlisting}[language=SQL]
CREATE OR REPLACE PROCEDURE proc_Emergency_Override (
    p_Intersection_ID   IN NUMBER,
    p_Emergency_Lane_ID IN NUMBER DEFAULT NULL
) IS
    v_Target_Lane NUMBER;
BEGIN
    IF p_Emergency_Lane_ID IS NOT NULL THEN
        v_Target_Lane := p_Emergency_Lane_ID;
    ELSE
        SELECT Lane_ID INTO v_Target_Lane FROM Lane
        WHERE Intersection_ID = p_Intersection_ID AND Is_Emergency_Corridor = 1
        FETCH FIRST 1 ROW ONLY;
    END IF;

    UPDATE Signal_Phase SET Current_State = 'RED', Duration_Seconds = 90
    WHERE Intersection_ID = p_Intersection_ID AND Lane_ID <> v_Target_Lane;

    UPDATE Signal_Phase SET Current_State = 'GREEN', Duration_Seconds = 60
    WHERE Intersection_ID = p_Intersection_ID AND Lane_ID = v_Target_Lane;
    
    COMMIT;
END proc_Emergency_Override;
/
\end{lstlisting}

\section{Reactive Shadow Audit Trigger}
\begin{lstlisting}[language=SQL]
CREATE OR REPLACE TRIGGER trg_Signal_Phase_Audit
AFTER UPDATE OF Current_State ON Signal_Phase
FOR EACH ROW
WHEN (OLD.Current_State IS NULL OR OLD.Current_State <> NEW.Current_State)
BEGIN
    INSERT INTO Audit_Traffic_Log (
        Audit_ID, Phase_ID, Old_State, New_State, Changed_At, Trigger_Source
    ) VALUES (
        seq_audit_traffic_log.NEXTVAL,
        :NEW.Phase_ID,
        :OLD.Current_State,
        :NEW.Current_State,
        SYSTIMESTAMP,
        'TRG_SIGNAL_PHASE_AUDIT'
    );
END;
/
\end{lstlisting}

% -----------------------------------------------------------------------------
% Chapter 6
% -----------------------------------------------------------------------------
\chapter{Experimental Verification \& Performance Analysis}

\section{Automated Verification Gates}
The system was validated using a 10-gate automated test suite (\texttt{scripts/test\_runner.py}):
\begin{itemize}
    \item \textbf{Gates 1--4:} Validated directory structure, DDL constraints, negative constraint violations, and 145+ seed records.
    \item \textbf{Gates 5--7:} Executed group analytics, 3-way density joins, subqueries, and verified that \texttt{vw\_Dashboard\_Live} isolated green phases with 100\% precision.
    \item \textbf{Gates 8--10:} Asserted atomic emergency overrides, verified weighted moving average bounds, and confirmed autonomous shadow audit insertions.
\end{itemize}

\section{Audit Log Verification Trace}
\begin{table}[ht]
\centering
\begin{tabular}{cccccc}
\toprule
\textbf{Audit ID} & \textbf{Phase ID} & \textbf{Old State} & \textbf{New State} & \textbf{Timestamp} & \textbf{Source} \\
\midrule
1 & 1 & GREEN & YELLOW & 2026-10-24 15:20:00 & TRG\_AUDIT \\
2 & 1 & YELLOW & RED & 2026-10-24 15:20:05 & TRG\_AUDIT \\
3 & 3 & RED & GREEN & 2026-10-24 15:20:10 & TRG\_AUDIT \\
\bottomrule
\end{tabular}
\caption{Verified Autonomous Audit Log Trail}
\end{table}

\section{Conclusion}
The SmartFlow Grid DBMS case study achieves complete adherence to relational theory and practical systems engineering. By combining rigorous BCNF normalization proofs with dual Oracle and PostgreSQL procedural logic, the architecture provides a robust, fail-safe backend for next-generation smart city traffic grids.

\end{document}
